% Partie : ensembles, fonctions et relations
% Chapitre : fonctions
% Section : isomorphismes

\documentclass[../../../include/open-logic-section]{subfiles}

\begin{document}

\olfileid[fr]{sfr}{fun}{iso}
\olsection{Isomorphismes}

\begin{explain}
Un \emph{isomorphisme} est une bijection qui
préserve la structure spécifiée sur les ensembles
qu'elle relie, ici les relations entre leurs
!!{element}s. Considérons
$X=\{1,2,3\}$ et $Y=\{4,5,6\}$,
munis des relations usuelles de successeur,
d'infériorité et de supériorité, restreintes
à chacun de ces ensembles. Une !!{bijection}
$f\colon X\to Y$ préserve ces trois relations
si, pour tous $i,j\in X$, on a
$i<j$ si et seulement si $f(i)<f(j)$,
$i>j$ si et seulement si $f(i)>f(j)$,
et $j$ est le successeur de $i$ si et seulement
si $f(j)$ est le successeur de $f(i)$.
Une telle bijection est un isomorphisme
pour ces structures munies des trois relations.
\end{explain}

\begin{defn}[Isomorphisme]
Soient $U=\langle X,R\rangle$ et
$V=\langle Y,S\rangle$, où $X$ et $Y$
sont des ensembles et $R$ et $S$ des
relations binaires respectivement sur $X$ et
sur $Y$. Une !!{bijection}~$f$ de $X$
sur $Y$ est un \emph{isomorphisme} de $U$
sur $V$ si elle préserve la relation
distinguée : pour tous $x_1,x_2\in X$,
$\tuple{x_1,x_2}\in R$ si et seulement si
$\tuple{f(x_1),f(x_2)}\in S$.
\end{defn}

\begin{ex}
Reprenons $X=\{1,2,3\}$ et
$Y=\{4,5,6\}$, mais distinguons
cette fois une seule relation dans chaque
structure : l'ordre strict $<$ sur~$X$
et l'ordre strict $>$ sur~$Y$.
La fonction $f\colon X\to Y$ définie
par $f(x)=7-x$ est un isomorphisme
de $\tuple{X,<}$ sur $\tuple{Y,>}$ :
pour $i,j\in X$, $i<j$ équivaut à
$f(i)>f(j)$, puisque $7-i>7-j$ équivaut à $i<j$.
L'image du successeur de~$1$ illustre la différence
entre les deux exemples. On a $f(2)=5$, qui est
le prédécesseur de~$f(1)=6$.
\end{ex}

\end{document}
